\scanpage{87}
\mspage{83}
% archive sheet number: 86

(cf construction de l'application $\unsure{\beta|B}$ de prop. 2), et au sens inverse par (iii) (application $\unsure{\alpha|\varphi}$ de prop. 2). Donc, (i) (ou ses voisines (iv), (v)) implique (ii), (iii) sous forme plus forte (iso. de \unsure{spotos|spatos}). Il reste à prouver que (ii) ou (iii) implique \inserted{déjà} (i).

Mais que (iii) (\unsure{et|i.e.} simplement la \uline{surjectivité} de $\psi$) implique (i), est clair. Donc (iii) équivaut aux autres conditions. Reste à voir que (ii) implique (i). Cela équivaut au corollaire suivant, appliqué à $E'\mathbin{-}\Sigma_E$

\uline{Corollaire 2} Soient $E$ un \unsure{disjoncteur}, $\Sigma$ un spotos, $\sigma:E\longrightarrow\Sigma$ une application. Alors la condition b) de la prop 2 est équivalente (moyennant a)) à la condition suivante:
\margin{\figure{\begin{tikzpicture}[x=1mm,y=1mm]\draw[->] (0,28)--(0,0);\end{tikzpicture}}}

b') L'application $S\longrightarrow\beta(S)$ de $\Sigma$ dans $\unsure{\mathscr{P}|\mathcal{P}|\mathfrak{P}}(E)$ applique $\Sigma$ dans $\Sigma_E$, et est \deleted{bijective}\inserted{injective}.
\margin{\figure{\begin{tikzpicture}[x=1mm,y=1mm]\draw (0,12)--(0,0);\end{tikzpicture}}}

On a vu dans la démonstration que b) $\Longrightarrow$ b'). Prouvons que b') $\Longrightarrow$ b). Considérons, pour $S\in\Sigma$ fixé, $S'=\operatorname{Sup}\sigma(\beta(S))\unsure{(}\leq S\unsure{)}$. Il est clair que $\beta(S)=\beta(S')$, donc $S=S'$, puisque $\beta$ est injectif qed
